\documentclass[10pt]{article}

\usepackage[utf8]{inputenc}

\usepackage[T1]{fontenc}

\usepackage{amsmath}

\usepackage{amsfonts}

\usepackage{amssymb}

\usepackage[version=4]{mhchem}

\usepackage{stmaryrd}

\usepackage{graphicx}

\usepackage[export]{adjustbox}

\graphicspath{ {./images/} }



\begin{document}

щ и м, или с ф е р ич е ским, или поч т и ин в ар и а н т пы м, если орбита $\varphi(G) v$ порождает в $E$ конечномерное подпространство. Это, в частности, применимо к случаю, когда $E$ - пространство сечений нек-рого класса гладкости гладкого векторного $G$ расслоения, напр. пространство тензорных полей определенного типа и заданного класса гладкости на гладком многообразии с гладким действием компактной групшы Ли $G$.



П.- В. т. была доказана в 1927 Ф. Петером (F. Peter) и Г. Вейлем (H. Weyl) (см. [1]).



गит.: [1] П е т е p Ф., В е й л b $\Gamma$., "Успехи матем. наук», 1936, в. 2, с. 144-60, [2] П о н т р я г и н J. С., Непрерывные рактный гармонический анализ, пер. с англ., т. 1, М., 1975; [4] III е в а л л е К., Теория групп Ји, пер. с англ., т. 1, М., $1948 ;$ [5] P al a i s R. S., S t e w a r t T. E.. "Amer. J. Math.", 1961,



\includegraphics[max width=\textwidth, center]{2023_07_08_c662f8e3f136787737e9g-1}

ПЕТЕРСОНА ПОВЕРХНОСТЬ - поверхность, обладающая сопряженной сетью конических или цилиндрич. линий, являющейся главным основанием изгибания (см. Изаибание на главном основании). Напр., П. П. является каналовая поверхность, переноса поверхность, вращения поверхность. Вращений индикатриса П. п. есть прямой коноид (в частности, для каналовой поверхности ею является геликоид, для поверхности переноса - гиперболич. параболоид). Впервые рассмотрена К. М. Петерсоном как пример поверхности, допускающей изгибание на главном основании.



IЕTEPСОНА СОOTBETCTBИЕ - соответствие двух поверхностей, при к-ром их касательные плоскости в соответствующих точках параллельны. В общем виде рассмотрено К. М. Петерсоном [1] в связи с задачей изгибания на главном основании. Напр., в П. с. находятся поверхность и ее сферич. образ, поверхность и ее индикатриса вращения, присоединенные минимальные поверхности.



Если поверхности $S$ и $S^{*}$ имеют общую параметризацию, то их третьи квадратичные формы равны. Главная сеть асимптотич. сетей $S$ и $S^{*}$ сопряжена на каждой из них. Эта сеть определяется однозначно, если асимптотич. сети не имеют общих семейств линий; вырождается, если эти сети связанные; становится неопределенной, если эти сети соответствуют друг другу. Соответствующие касательные к линиям главной сети на $S$ и $S^{*}$ параллельны.



Если главную сеть принять за координатную $(u, v)$, то радиус-векторы $x$ и $x^{*}$ поверхностей $S$ и $S^{*}$ связаны уравнениями



\[

x_{u}^{*}=\rho x_{u}, \quad x_{v}^{*}=\sigma x_{v},

\]



причем функции $\rho, \sigma$ удовлетворяют системе уравнени й



\[

\rho_{v} E-\sigma_{u} F=\frac{\sigma-\rho}{2}, \rho_{v} F-\sigma_{u} G=\frac{\sigma-\rho}{2} G_{u},

\]



то есть $\rho, \sigma$ зависят только от метрики $F(E, F, G-$ коэффициенты ее первой квадратичной формы). Естественно поэтому применение П. с. к паре изометричных поверхностей $S, \tilde{S}$ : оно дает другую пару изометричных поверхностей $S^{*}, \tilde{S^{*}}$ с теми же нормалями $n=n^{*}, \tilde{n}=\tilde{n}$ соответственно. Кроме того, оказывается, что диаграммы поворота этих поверхностей одни и те же и что основание изометрии новых поверхностей имеет то же сферич. изображение, как и первоначальных. Напр., сфере и изометричной ей поверхности положительной постоянной кривизны отвечают при П. с. изометричные поверхности с соответствующими линиями кривизны - т. н. поверхности Бонне.



В частности, если основание изометрии $S$ и $S^{*}$ было главным основанием, то оно и остается таковым. Это-т. н. преобра зование Пете рсона поверхности, изгибаемой на главном основании в поверхность того же типа. Имеется обобщение этого преобразования на случай семейства сетей (см. [2]).



Специальный случай П. с., при к-ром главная сеть является сетью кривизны одновременно на $S$ и на $S^{*}$, наз. с о о т в е т с т и е м $К$ о м бес к ю р а.



Если П. с. является конформным, то либо одна из поверхностей минимальная, а другая - сфера (то есть П. с. - сферич. отображение), либо обе поверхности минимальные, а вектор конформного отображения удовлетворяет уравнению Лапласа, либо поверхности подобны, либо обе поверхности изотермические и находятся в с о о т в е т т в И и К о б е ск ю р а.



Имеется многомерное обобщение П. с. (см. [4]).



Лum.: [1] П е т е р с о в К. М., "Матем. сб.", 1866, т. 1, с.

$391-438 ;$ [2] Ш у л и к о в с к и й В. И., Классическая дифференциальная геометрия в тензорном изложении, М., 1963; [3] Ф и н и к о в С. П., Изгибание на главном основании и связанные с ним геометрические задачи, М.-- J., 1937; [4] IШ и р ок о в П. А., ІІ и ро к о в А. П., Аффинная дифференцияьй

геометрия, М., 1959.



ПЕТЕРСОНА - КОДАЦЦИ УРАВНЕНИЯ - уравнения, составляющие вместе с уравнением Гаусса (см. Гаусса теорема) необходимые и достаточные условия интегрируемости системы, к к-рой сводится задача восстановления поверхности по ее первой и второй квадратичным формам (см. Бонне теорема). П. - К. у. имеют вид



\[

\frac{\partial b_{i 1}}{\partial u^{2}}+\Gamma_{i 1}^{1} b_{12}+\Gamma_{i 1}^{2} b_{22}=\frac{\partial b_{i 2}}{\partial u^{1}}+\Gamma_{i 2}^{1} b_{11}+\Gamma_{i 2}^{2} b_{21},

\]



где $b_{i j}$ - коэффициенты второй квадратичной формы, $\Gamma_{i j}^{k}$ - символы Кристоффеля 2-го рода.



впервые найдены К. Петерсоном в 1853, переоткрыты Г. Майнарди (G. Mainardi, 1856) и Д. Кодацци (D. Codazzi, 1867).



Лum.: [1] Р a шा е в с к и й П. К., Курс дифференциальной геометрии, М., 1956. А. Б. Иванов.



ІЕТЛЯ - замкнутый путь. Подробнее, петля $f$ непрерывное отображение отрезка $[0,1]$ в топологич. пространство $X$ такое, что $f(0)=f(1)$. Совокупность всех П. пространства $\dot{X}$ с отмеченной точкой * таких, что $f(0)=f(1)=*$, образует петель пространство.



IЕТРИ СЕТЬ - М. И. Войцеховский. динамич. систем, в том числе информационных систем (параллельных программ, операционных систем, ЭВМ и их устройств, сетей ЭВМ), ориентированная на качественный анализ и синтез таких систем (обнаружение блокировок, тупиковых ситуаций и узких мест, автоматич: синтез параллельных программ и компонентов ЭВМ и др.). Введена К. Петри (C. Petri) в 60-х гг. 20 в. П. с. - это набор $N=\left(T, P, F, M_{0}\right)$, где $T-$ конечное множество символов, наз. П е $\mathrm{p} \mathrm{e} \mathrm{x} о$ д а м и, $P$ - конечное множество символов, наз. м е с т а м и, $P \cap T=\phi, F$ - функция инцидентности:



\[

F: T \times P \cup P \times T \longrightarrow\{0,1\}

\]



$M_{0}$ - начальная разметка мест:



\[

M_{0}: P \longrightarrow\{0,1,2, \ldots\}

\]



Неформально П. с. представляют размеченным ориентированным графом со множеством вершпин $T \cup P$ (см. рис.). Из вершины-места $p \in P$, изображаемой кружком, ведет дуга в вершину-переход $t \in T$, изображаемую прямоугольником, тогда и только тогда, когда



\[

F(p, t)=1

\]



( $p$-входное место перехода $t$; на рис. $P=\left\{p_{1}, p_{2}, p_{3}\right\}$, $T=\{a, b, c, d\})$. Из вершины-перехода $t$ ведет дуга в вершину-место $p$ тогда џ только тогда, когда



\[

F(t, p)=1

\]





\end{document}